\studySubsection{practice-calc-09}{练习库：高等数学第 9 讲 一元函数积分学的计算}

\problemAnchor{MATH1-CALC-0031}
\begin{problemBox}
\begin{problemMeta}
\textbf{题目编号：} \problemRef{MATH1-CALC-0031}\quad
\textbf{答案：} \answerRef{MATH1-CALC-0031}\quad
\textbf{来源：} Codex 原创 / K082 深度讲义配套练习（非真题）
\end{problemMeta}
\begin{enumerate}
  \item[\textbf{A}] 求 \(\displaystyle\int3x^2\cos(x^3)\,\mathrm dx\)。
  \item[\textbf{B}] 求 \(\displaystyle\int\frac{x}{(1+x^2)^2}\,\mathrm dx\)。
  \item[\textbf{C}] 在 \(x>0\) 上求
  \(\displaystyle\int\frac{e^{\sqrt x}}{\sqrt x}\,\mathrm dx\)。
  \item[\textbf{M}] 设 \(\lambda\in\mathbb R\)。在
  \(x^2+\lambda x+1\ne0\) 的任一连通区间上求
  \[
  \int\frac{2x+\lambda}{x^2+\lambda x+1}\,\mathrm dx.
  \]
\end{enumerate}
\end{problemBox}

\problemAnchor{MATH1-CALC-0032}
\begin{problemBox}
\begin{problemMeta}
\textbf{题目编号：} \problemRef{MATH1-CALC-0032}\quad
\textbf{答案：} \answerRef{MATH1-CALC-0032}\quad
\textbf{来源：} Codex 原创 / K083 深度讲义配套练习（非真题）
\end{problemMeta}
\begin{enumerate}
  \item[\textbf{A}] 求
  \(\displaystyle\int\frac{\mathrm dx}{\sqrt{9-x^2}}\)。
  \item[\textbf{B}] 在 \(|x|<2\) 上求
  \(\displaystyle\int\frac{x^2}{\sqrt{4-x^2}}\,\mathrm dx\)。
  \item[\textbf{C}] 求
  \(\displaystyle\int\frac{\mathrm dx}{(x^2+4)^{3/2}}\)。
  \item[\textbf{M}] 设 \(a>0\)，不用几何面积公式，计算
  \[
  \int_0^a x^2\sqrt{a^2-x^2}\,\mathrm dx.
  \]
\end{enumerate}
\end{problemBox}

\problemAnchor{MATH1-CALC-0033}
\begin{problemBox}
\begin{problemMeta}
\textbf{题目编号：} \problemRef{MATH1-CALC-0033}\quad
\textbf{答案：} \answerRef{MATH1-CALC-0033}\quad
\textbf{来源：} Codex 原创 / K084 深度讲义配套练习（非真题）
\end{problemMeta}
\begin{enumerate}
  \item[\textbf{A}] 求 \(\displaystyle\int x\cos x\,\mathrm dx\)。
  \item[\textbf{B}] 求 \(\displaystyle\int\arctan x\,\mathrm dx\)。
  \item[\textbf{C}] 求 \(\displaystyle\int x^2e^{2x}\,\mathrm dx\)。
  \item[\textbf{M}] 对整数 \(n\ge2\)，令
  \[
  J_n=\int_0^{\pi/2}\sin^n x\,\mathrm dx.
  \]
  用分部积分证明
  \(J_n=\frac{n-1}{n}J_{n-2}\)，并给出 \(J_0,J_1\)。
\end{enumerate}
\end{problemBox}

\problemAnchor{MATH1-CALC-0034}
\begin{problemBox}
\begin{problemMeta}
\textbf{题目编号：} \problemRef{MATH1-CALC-0034}\quad
\textbf{答案：} \answerRef{MATH1-CALC-0034}\quad
\textbf{来源：} Codex 原创 / K085 深度讲义配套练习（非真题）
\end{problemMeta}
\begin{enumerate}
  \item[\textbf{A}] 求
  \(\displaystyle\int\frac{2x+1}{x^2+x}\,\mathrm dx\)。
  \item[\textbf{B}] 求
  \(\displaystyle\int\frac{x^2+2}{x^2-1}\,\mathrm dx\)。
  \item[\textbf{C}] 求
  \[
  \int\frac{x^2+x+1}{(x-1)^2(x+1)}\,\mathrm dx.
  \]
  \item[\textbf{M}] 求
  \[
  \int\frac{x^2+1}{x(x^2+4)}\,\mathrm dx.
  \]
\end{enumerate}
\end{problemBox}

\problemAnchor{MATH1-CALC-0035}
\begin{problemBox}
\begin{problemMeta}
\textbf{题目编号：} \problemRef{MATH1-CALC-0035}\quad
\textbf{答案：} \answerRef{MATH1-CALC-0035}\quad
\textbf{来源：} Codex 原创 / K086 深度讲义配套练习（非真题）
\end{problemMeta}
\begin{enumerate}
  \item[\textbf{A}] 求 \(\displaystyle\int\sin^3x\,\mathrm dx\)。
  \item[\textbf{B}] 求 \(\displaystyle\int\cos^4x\,\mathrm dx\)。
  \item[\textbf{C}] 在被积函数有定义的连通区间上求
  \[
  \int\frac{\mathrm dx}{\sin x\cos x}.
  \]
  \item[\textbf{M}] 使用万能代换求
  \[
  \int\frac{\mathrm dx}{2+\sin x+\cos x}.
  \]
\end{enumerate}
\end{problemBox}

\problemAnchor{MATH1-CALC-0036}
\begin{problemBox}
\begin{problemMeta}
\textbf{题目编号：} \problemRef{MATH1-CALC-0036}\quad
\textbf{答案：} \answerRef{MATH1-CALC-0036}\quad
\textbf{来源：} Codex 原创 / K087 深度讲义配套练习（非真题）
\end{problemMeta}
\begin{enumerate}
  \item[\textbf{A}] 在 \(x>0\) 上求
  \(\displaystyle\int\frac{x}{1+\sqrt x}\,\mathrm dx\)。
  \item[\textbf{B}] 在 \(x>0\) 上求
  \(\displaystyle\int\frac{\mathrm dx}{\sqrt x(1+\sqrt x)}\)。
  \item[\textbf{C}] 在 \(x>0\) 上求
  \(\displaystyle\int\frac{\sqrt x}{1+x}\,\mathrm dx\)。
  \item[\textbf{M}] 在 \(x>-1\) 且
  \(x+\sqrt{x+1}\ne0\) 的任一连通区间上求
  \[
  \int\frac{\mathrm dx}{x+\sqrt{x+1}}.
  \]
\end{enumerate}
\end{problemBox}

\problemAnchor{MATH1-CALC-0042}
\begin{problemBox}
\begin{problemMeta}
\textbf{题目编号：} \problemRef{MATH1-CALC-0042}\quad
\textbf{答案：} \answerRef{MATH1-CALC-0042}\quad
\textbf{来源：} Codex 原创 / K093 深度讲义配套练习（非真题）
\end{problemMeta}
\begin{enumerate}
  \item[\textbf{A}] 计算
  \(\displaystyle\int_0^1\frac{3x^2}{1+x^3}\,\mathrm dx\)。
  \item[\textbf{B}] 计算
  \(\displaystyle\int_0^2x\sqrt{4-x^2}\,\mathrm dx\)。
  \item[\textbf{C}] 设 \(a>0\)，\(f\) 在 \([0,a]\) 可积。证明
  \[
  \int_0^a x\bigl(f(x)+f(a-x)\bigr)\,\mathrm dx
  =a\int_0^a f(x)\,\mathrm dx.
  \]
  \item[\textbf{M}] 计算
  \[
  \int_0^{2\pi}|\sin x|e^{\cos x}\,\mathrm dx,
  \]
  并说明为什么不能在整个区间上直接令 \(u=\cos x\) 后只写一组
  新上下限。
\end{enumerate}
\end{problemBox}

\problemAnchor{MATH1-CALC-0043}
\begin{problemBox}
\begin{problemMeta}
\textbf{题目编号：} \problemRef{MATH1-CALC-0043}\quad
\textbf{答案：} \answerRef{MATH1-CALC-0043}\quad
\textbf{来源：} Codex 原创 / K094 深度讲义配套练习（非真题）
\end{problemMeta}
\begin{enumerate}
  \item[\textbf{A}] 计算反常积分
  \(\displaystyle\int_0^1x\ln x\,\mathrm dx\)。
  \item[\textbf{B}] 计算
  \(\displaystyle\int_0^1x\arctan x\,\mathrm dx\)。
  \item[\textbf{C}] 对 \(a>-1\)，计算
  \[
  I(a)=\int_0^1x^a(\ln x)^2\,\mathrm dx.
  \]
  \item[\textbf{M}] 先在有限区间分部，再计算
  \[
  \int_1^\infty\frac{\ln x}{x^2}\,\mathrm dx.
  \]
\end{enumerate}
\end{problemBox}

\problemAnchor{MATH1-CALC-0044}
\begin{problemBox}
\begin{problemMeta}
\textbf{题目编号：} \problemRef{MATH1-CALC-0044}\quad
\textbf{答案：} \answerRef{MATH1-CALC-0044}\quad
\textbf{来源：} Codex 原创 / K095 深度讲义配套练习（非真题）
\end{problemMeta}
\begin{enumerate}
  \item[\textbf{A}] 计算
  \[
  \int_0^2(x-1)^7e^{(x-1)^2}\,\mathrm dx.
  \]
  \item[\textbf{B}] 设 \(f\) 在 \([0,a]\) 可积并满足
  \(f(a-x)=f(x)\)。证明
  \[
  \int_0^a\left(x-\frac a2\right)^3f(x)\,\mathrm dx=0.
  \]
  \item[\textbf{C}] 设 \(f\) 以 \(T>0\) 为周期且局部可积。对正整数
  \(n\)，证明
  \[
  \int_c^{c+nT}f(x)\,\mathrm dx
  =n\int_0^Tf(x)\,\mathrm dx.
  \]
  \item[\textbf{M}] 计算
  \[
  \int_0^{2\pi}\frac{x}{2+\cos x}\,\mathrm dx.
  \]
\end{enumerate}
\end{problemBox}

\problemAnchor{MATH1-CALC-0050}
\begin{problemBox}
\begin{problemMeta}
\textbf{题目编号：} \problemRef{MATH1-CALC-0050}\quad
\textbf{答案：} \answerRef{MATH1-CALC-0050}\quad
\textbf{来源：} Codex 原创 / K101 深度讲义配套练习（非真题）
\end{problemMeta}
\begin{enumerate}
  \item[\textbf{A}] 计算
  \(\displaystyle\int_0^\infty xe^{-x^2}\,\mathrm dx\)。
  \item[\textbf{B}] 计算
  \(\displaystyle\int_{-\infty}^{\infty}
  \frac{e^x}{(1+e^x)^2}\,\mathrm dx\)。
  \item[\textbf{C}] 讨论参数 \(a\) 取何值时
  \[
  \int_0^\infty xe^{-ax^2}\,\mathrm dx
  \]
  收敛，并在收敛时求值。
  \item[\textbf{M}] 证明
  \[
  \int_{-\infty}^{\infty}\frac{x}{1+x^2}\,\mathrm dx
  \]
  作为通常反常积分发散，但其 Cauchy 主值为 \(0\)。
\end{enumerate}
\end{problemBox}

\problemAnchor{MATH1-CALC-0051}
\begin{problemBox}
\begin{problemMeta}
\textbf{题目编号：} \problemRef{MATH1-CALC-0051}\quad
\textbf{答案：} \answerRef{MATH1-CALC-0051}\quad
\textbf{来源：} Codex 原创 / K102 深度讲义配套练习（非真题）
\end{problemMeta}
\begin{enumerate}
  \item[\textbf{A}] 计算
  \(\displaystyle\int_0^1\frac{\mathrm dx}{\sqrt x+x}\)。
  \item[\textbf{B}] 计算
  \(\displaystyle\int_0^1\ln x\,\mathrm dx\)。
  \item[\textbf{C}] 讨论
  \[
  \int_0^2\frac{\mathrm dx}{|x-1|^p}
  \]
  的收敛性，并在收敛时求值。
  \item[\textbf{M}] 对 \(c\in\mathbb R\)，讨论
  \[
  \int_{-1}^{1}\frac{\mathrm dx}{x-c}
  \]
  何时为普通定积分或收敛反常积分、何时发散；在收敛情形求值。
\end{enumerate}
\end{problemBox}

\problemAnchor{MATH1-CALC-0052}
\begin{problemBox}
\begin{problemMeta}
\textbf{题目编号：} \problemRef{MATH1-CALC-0052}\quad
\textbf{答案：} \answerRef{MATH1-CALC-0052}\quad
\textbf{来源：} Codex 原创 / K103 深度讲义配套练习（非真题）
\end{problemMeta}
\begin{enumerate}
  \item[\textbf{A}] 计算
  \(\displaystyle\int_1^\infty\frac{\mathrm dx}{x^{3/2}+x}\)。
  \item[\textbf{B}] 计算
  \(\displaystyle\int_0^1
  \frac{\mathrm dx}{x^{2/3}+x^{1/3}}\)。
  \item[\textbf{C}] 讨论参数 \(\alpha,\beta\) 取何值时
  \[
  \int_0^{\pi/2}(\sin x)^\alpha(\cos x)^\beta\,\mathrm dx
  \]
  收敛。
  \item[\textbf{M}] 讨论
  \[
  \int_{e^e}^\infty
  \frac{\mathrm dx}{x(\ln x)^p(\ln\ln x)^q}
  \]
  随参数 \(p,q\) 的收敛性。
\end{enumerate}
\end{problemBox}

\problemAnchor{MATH1-CALC-0053}
\begin{problemBox}
\begin{problemMeta}
\textbf{题目编号：} \problemRef{MATH1-CALC-0053}\quad
\textbf{答案：} \answerRef{MATH1-CALC-0053}\quad
\textbf{来源：} Codex 原创 / K104 深度讲义配套练习（非真题）
\end{problemMeta}
\begin{enumerate}
  \item[\textbf{A}] 判断
  \(\displaystyle\int_1^\infty\frac{\mathrm dx}{x^2+\sqrt x}\)
  的敛散性。
  \item[\textbf{B}] 判断
  \(\displaystyle\int_0^1\frac{\mathrm dx}{\sqrt x(1+x)}\)
  的敛散性。
  \item[\textbf{C}] 讨论
  \[
  \int_0^1\frac{1-\cos x}{x^p}\,\mathrm dx
  \]
  随参数 \(p\) 的敛散性。
  \item[\textbf{M}] 讨论
  \[
  \int_1^\infty
  \frac{\sqrt{x^2+x}-x}{x^p}\,\mathrm dx
  \]
  随参数 \(p\) 的敛散性。
\end{enumerate}
\end{problemBox}

\problemAnchor{MATH1-CALC-0054}
\begin{problemBox}
\begin{problemMeta}
\textbf{题目编号：} \problemRef{MATH1-CALC-0054}\quad
\textbf{答案：} \answerRef{MATH1-CALC-0054}\quad
\textbf{来源：} Codex 原创 / K105 深度讲义配套练习（非真题）
\end{problemMeta}
\begin{enumerate}
  \item[\textbf{A}] 证明
  \(\displaystyle\int_0^\infty e^{-x}\sin x\,\mathrm dx\)
  绝对收敛。
  \item[\textbf{B}] 证明
  \(\displaystyle\int_1^\infty\frac{\cos x}{\sqrt x}\,\mathrm dx\)
  条件收敛。
  \item[\textbf{C}] 判断
  \(\displaystyle\int_1^\infty\sin(x^2)\,\mathrm dx\)
  是绝对收敛、条件收敛还是发散。
  \item[\textbf{M}] 讨论
  \[
  \int_1^\infty
  \frac{\sin x}{x^p(1+x)}\,\mathrm dx
  \]
  随参数 \(p\) 的绝对收敛、条件收敛与发散情形。
\end{enumerate}
\end{problemBox}

\problemAnchor{MATH1-CALC-0061}
\begin{problemBox}
\begin{problemMeta}
\textbf{题目编号：} \problemRef{MATH1-CALC-0061}\quad
\textbf{答案：} \answerRef{MATH1-CALC-0061}\quad
\textbf{来源：} Codex 原创 / K112 深度讲义配套练习（非真题）
\end{problemMeta}
\begin{enumerate}
  \item[\textbf{A}] 不使用特殊函数，计算
  \(\displaystyle\int_0^\infty e^{-\sqrt x}\,\mathrm dx\)。
  \item[\textbf{B}] 先截断再分部，计算
  \(\displaystyle\int_0^\infty\frac{x}{(1+x)^3}\,\mathrm dx\)。
  \item[\textbf{C}] 不使用 Beta 函数公式，计算
  \[
  \int_0^1\sqrt{x(1-x)}\,\mathrm dx.
  \]
  \item[\textbf{M}] 对非负整数 \(n\)，仅用反常积分定义与分部积分证明
  \[
  \int_0^\infty x^ne^{-x}\,\mathrm dx=n!.
  \]
\end{enumerate}
\end{problemBox}
